\documentclass{article}
\usepackage[T1]{fontenc}
\usepackage{lmodern}
\usepackage{graphicx}
\usepackage{amsmath,amssymb}
\usepackage[a4paper,margin=2cm]{geometry}
\usepackage[absolute,overlay]{textpos}
\setlength{\TPHorizModule}{1mm}
\setlength{\TPVertModule}{1mm}
\usepackage[indonesian]{babel}
\usepackage{hyperref}
\setlength{\emergencystretch}{2em}
% unit: Halaman 15

\begin{document}

% === Halaman 15 (python/native) ===

15

• Bob mengenkripsi setiap blok plainteks dengan menggunakan kunci publik 

Alice (e = 79, n = 3337):


m1 = 0704 →c1 = 70479 mod 3337 = 328;

m2 = 1111 →c2 = 111179 mod 3337 = 301;

m3 = 1400 →c3 = 140079 mod 3337 = 2653;

m4 = 1108 →c4 = 110879 mod 3337 = 2986; 

m1 = 204 →c5 = 20479 mod 3337 = 1164; 


    Cipherteks: C = 0328, 0301, 2653, 2986, 1164

.

• Bob mengirim cipherteks C kepada Alice

Enkripsi: c = me mod n

\end{document}
